\originalchapter{8.334作业2：质量生成、随机场与长程相互作用}
\psource{Mehran Kardar，8.334 Spring 2014，PS2，截止2014年3月4日。}
\begin{exercise}[原题1(a)--(j)]\label{p:334-2-1}
$n$分量$m$与标量$A$耦合，$\beta H=\int[K(\nabla m)^2/2+tm^2/2+u(m^2)^2+em^2A^2+L(\nabla A)^2/2]\dd^dx$，$K,L,u>0$，为稳定取$e\ge0$。(a)均匀鞍点$m\hat e_\ell,A=0$；(b)$C=\partial_t^2\log Z$；(c)分纵横与$a$模二阶展开；(d)纵横相关长；(e)$a$相关长；(f)$t>0$时$\avg{a(x)a(0)}$；(g)Gaussian自由能修正；(h)热容修正；(i)积分随$\xi,\Lambda,d$行为；(j)鞍点上临界维是否由标量耦合改变。
\end{exercise}
\begin{solution}
$m_0=0$（$t>0$）或$m_0^2=-t/(4u)$（$t<0$），鞍点$f_0=0$或$-t^2/(16u)$。所以$C_0/V=0$或$1/(8u)$，为阶跃。置$m=(m_0+\phi_\ell)e_\ell+\sum\phi_te_t$、$A=a$，Hessian质量$r_\ell=t+12um_0^2$、$r_t=t+4um_0^2$、$r_a=2em_0^2$；没有二次混合项，因为原耦合最少$a^2$。

纵横相关长$\sqrt{K/r_{\ell,t}}$：正侧皆$\sqrt{K/t}$，负侧纵$\sqrt{K/(-2t)}$、横无穷（Goldstone）；$\xi_a$正侧无穷，负侧$\sqrt{2uL/(-et)}$。正侧$G_a(r)=\int_q\e^{\ii qx}/(Lq^2)$，$d>2$远距$\Gamma(d/2-1)/(4\pi^{d/2}Lr^{d-2})$，$d\le2$无限体积绝对场相关需IR规整。

Gaussian积分得
\begin{equation}
\delta\log Z/V=-\tfrac12\int_q[\log(Kq^2+r_\ell)+(n-1)\log(Kq^2+r_t)+\log(Lq^2+r_a)],\quad \int_q=\int_{|q|<\Lambda}\frac{\dd^dq}{(2\pi)^d}.
\end{equation}
省去与$t$无关测度常数。各分支质量随$t$线性，$\delta C/V=\tfrac12\int_q\sum_i(r_i')^2/(K_iq^2+r_i)^2$。正侧为$n\int_q(Kq^2+t)^{-2}/2$；负侧为$[4\int_q(Kq^2-2t)^{-2}+(e/2u)^2\int_q(Lq^2-et/(2u))^{-2}]/2$。横模质量恒0，不产生这阶$t$导数；仅$t=0$不可用分段微分跨越鞍点。

$\int_q(q^2+\xi^{-2})^{-2}$在$d<4$奇异$\propto\xi^{4-d}$，$d=4$为$\log(\Lambda\xi)$，$d>4$UV主导$\Lambda^{d-4}$加较小奇异项，故Gaussian判据仍$d_u=4$。$A$增加奇异项系数，不改变这一维数；不过积分掉标量可生成非解析有效势，某些参数下使连续转变被一阶转变取代，该一阶问题不由二阶鞍点计算自动排除。
\end{solution}
\begin{exercise}[原题2(a)--(i)]\label{p:334-2-2}
标量$\beta H=\int[K(\nabla m)^2/2+tm^2/2+um^4-h(x)m]\dd^dx$，$u>0$，淬火Gaussian场$\overline h=0$、$\overline{h(x)h(y)}=\Delta\delta(x-y)$，等价$\overline{h(q)h(q')}=\Delta(2\pi)^d\delta(q+q')$。(a)均匀鞍点淬火平均；(b)二阶涨落；(c)Fourier能；(d)固定$h$下均值与连接方差；(e)平均散射线形；(f),(g)固定$h$及平均Gaussian自由能；(h)奇异$t$幂次；(i)上临界维。
\end{exercise}
\begin{solution}
均匀假设仅感受$\bar h_V=V^{-1}\int h$，其方差$\Delta/V$，热力学平均消失，故$m_0=0$或$m_0^2=-t/(4u)$、$f_{sp}=0$或$-t^2/(16u)$。但真实随机场局部变形不能由均匀平均排除。

置$m=m_0+\phi$，二次核$D(q)=Kq^2+r$、$r=t+12um_0^2$，$\beta H_2=\tfrac12\int_qD(q)|\phi(q)|^2-\int_q h(-q)\phi(q)$。配方$\phi=h/D+\delta\phi$，得到$\avg{\phi(q)}_h=h(q)/D$、$\avg{\delta\phi(q)\delta\phi(q')}_h=(2\pi)^d\delta(q+q')/D$。淬火后总线形
\begin{equation}
S(q)=D(q)^{-1}+\Delta D(q)^{-2},\quad
\delta f[h]=\tfrac12\int_q\log D-\frac1{2V}\int_q|h(q)|^2/D.
\end{equation}
第二项不是热连接方差；淬火平均$\overline{\delta f}=\tfrac12\int_q\log D-\tfrac\Delta2\int_qD^{-1}$。减去解析UV项后热奇异$r^{d/2}$，随机项$-\Delta r^{(d-2)/2}$，偶数边界出现对数。负側平均场$f_{sp}\sim t^2$；随机修正相对它$\sim\Delta |t|^{(d-6)/2}$，故上临界维6，不能由清洁系统Ginzburg的4代替。此Gaussian判据不等价于全非微扰随机场临界解，也不证明维数约化。
\end{solution}
\begin{exercise}[原题3(a)--(c)]\label{p:334-2-3}
连续自旋除短程项外有铁磁尾$J(r)\propto r^{-d-\sigma}$。(a)二次核的变化及何时主导；(b)由Goldstone涨落求下临界维；(c)鞍点上临界维。
\end{exercise}
\begin{solution}
差能量用$\int\dd^drJ(r)[1-\cos(qr)]$，$0<\sigma<2$作$r=y/q$得到$|q|^\sigma$，系数正；$\sigma>2$二阶矩有限，领先$q^2$，等号有$q^2\log q$。因此Gaussian层长程主导$\sigma<2$。Goldstone方差$\int_0^\Lambda q^{d-1}\dd q/q^\sigma$在$d\le\sigma$IR发散，所以连续对称$d_\ell=\sigma$（若短程主导则2）。场维$(d-\sigma)/2$使$um^4$尺度维$2\sigma-d$，$d_u=2\sigma$；相互作用临界下长短程精确交界可移到$\sigma=2-\eta_{SR}$，不能把Gaussian阈值当所有参数的全阶结论。
\end{solution}
